\section{Methods}\label{sec:methods}
\owner{Ky (K-Means, K-Means++, DBSCAN, HDBSCAN); Lam (GMM, Spectral); Hoang Anh (Agglomerative)}

% Rubric: "clear explanation of the tasks and the methods". Each method: idea, objective or
% rule, parameters, complexity, assumption it makes, how it fails.

\subsection{K-Means and K-Means++}
% Ky. Code: SC4020-Ky/src/core.py (run_kmeans)
K-means~\cite{lloyd1982} looks for $k$ centres $\boldsymbol{\mu}_1,\dots,\boldsymbol{\mu}_k$ that minimise
the within-cluster sum of squares (inertia), $J=\sum_{i=1}^{n}\min_j\lVert\mathbf{x}_i-\boldsymbol{\mu}_j\rVert^2$.
Each iteration assigns every point to its nearest centre and moves each centre to the mean of its points.
$J$ never rises, so the method stops at a local optimum that depends on the starting centres. One
iteration costs $O(nkd)$. Random starts pick $k$ data points uniformly. K-means++~\cite{arthur2007} picks
the first centre at random and each next one with probability proportional to its squared distance from
the nearest centre already chosen. Only the start differs, so we treat the two as one method with two
initialisations (Ablation~A). We use scikit-learn's \texttt{KMeans} with 10 starts; the only parameter is
$k$. Because $J$ sums squared Euclidean distances, K-means suits compact, roughly round clusters of
similar spread. It cuts curved or elongated clusters, gives every point a cluster, so outliers pull the
centres, and depends on the scale of each feature (Ablation~C).
\subsection{DBSCAN}
% Ky. Code: SC4020-Ky/src/core.py (run_dbscan, k_distance_curve, knee)
DBSCAN~\cite{ester1996} takes a radius $\varepsilon$ and a count \texttt{minPts}. A point is a core point
if at least \texttt{minPts} points lie within $\varepsilon$ of it. Core points within $\varepsilon$ of
each other join the same cluster, points within $\varepsilon$ of a core point join its cluster as border
points, and every other point is noise. The method needs no $k$ and assumes no shape, and on coordinates
$\varepsilon$ is a distance on the ground that an analyst can judge. With a spatial index a fit costs
about $O(n\log n)$ in few dimensions, rising towards $O(n^2)$ in many dimensions or with a large
$\varepsilon$. A reference value for $\varepsilon$ comes from the k-distance curve, each point's distance
to its \texttt{minPts}-th nearest neighbour in sorted order; its knee is the point farthest from the line
joining the two ends of the curve. The weakness is the single radius. It assumes all clusters have
similar density, so a radius that separates dense clusters calls sparse ones noise, and one that keeps
sparse clusters merges dense ones. In many dimensions distances become alike~\cite{beyer1999}, and a
radius separates little.
\subsection{HDBSCAN}
% Ky. Code: SC4020-Ky/src/core.py (run_hdbscan)
HDBSCAN~\cite{campello2013} removes $\varepsilon$. Each point's core distance is the distance to its
$m$-th nearest neighbour, and the mutual reachability distance
$d_{\mathrm{mr}}(a,b)=\max\{\mathrm{core}(a),\mathrm{core}(b),d(a,b)\}$ pushes points in sparse regions
apart. A minimum spanning tree under $d_{\mathrm{mr}}$ holds the DBSCAN clusterings at every density
level at once. The tree is condensed by dropping any split that leaves fewer than
\texttt{min\_cluster\_size} points, and the clusters that persist over the widest range of density are
kept; the remaining points are noise. Different branches can thus be cut at different densities, which is
what DBSCAN cannot do. We use scikit-learn's \texttt{HDBSCAN}; the main parameter is
\texttt{min\_cluster\_size}, and $m$ (\texttt{min\_samples}) defaults to it. A fit costs between
$O(n\log n)$ and $O(n^2)$. Because \texttt{min\_cluster\_size} is a count, the result changes with the
size of the sample. HDBSCAN can call large parts of the data noise, and like DBSCAN it has nothing to
find where the data have no density gaps.
\subsection{Gaussian Mixture Model}
% Lam. Code: lam/src/sc4020/clustering.py
A Gaussian mixture model (GMM) treats the data as drawn from $k$ Gaussians,
$p(\mathbf{x})=\sum_{j=1}^{k}\pi_j\,\mathcal{N}(\mathbf{x}\mid\boldsymbol{\mu}_j,\boldsymbol{\Sigma}_j)$,
and fits the weights $\pi_j$, means and covariances by expectation--maximisation (EM)~\cite{dempster1977}.
The E-step computes each point's responsibility (posterior probability) under every component; the
M-step re-estimates $\pi_j$, $\boldsymbol{\mu}_j$ and $\boldsymbol{\Sigma}_j$ from those soft
assignments; the two steps repeat until the log-likelihood stops improving. Each point is then given
its most probable component. We use scikit-learn's \texttt{GaussianMixture}~\cite{pedregosa2011} with
\emph{full} covariance matrices (the default), $k$-means initialisation, one start and a fixed seed.
The only parameter is the number of components $k$, which we set to the number of classes on labelled
data and to the silhouette-best $k\in[2,10]$ otherwise. One EM iteration costs $O(nkd^2+kd^3)$.
GMM generalises K-Means: full covariances allow ellipsoidal clusters of different size, orientation
and weight, but each cluster must still be a single convex blob. It fails on curved or nested shapes,
has no noise label, and with full covariances it needs many points per component when $d$ is large.
\subsection{Spectral clustering}
% Lam. Code: lam/src/sc4020/clustering.py; affinity per dataset in lam/notebooks/0*_*.ipynb
Spectral clustering~\cite{ng2002} turns clustering into graph partitioning. It (i) builds an affinity
matrix $W$ whose entry $w_{ij}$ measures how similar points $i$ and $j$ are; (ii) forms the normalised
graph Laplacian and takes its $k$ leading eigenvectors, giving each point a $k$-dimensional embedding;
(iii) runs K-Means on the rows of that embedding. Points joined by a chain of close neighbours land
near each other in the embedding even when they are far apart in the input space, so the method can
recover curved and non-convex clusters. We use scikit-learn's \texttt{SpectralClustering} with two
affinities: an RBF kernel $w_{ij}=\exp(-\gamma\lVert\mathbf{x}_i-\mathbf{x}_j\rVert^2)$, $\gamma=1$,
for roughly globular data (D6 Wine, Mall), and a symmetric 10-nearest-neighbour graph for manifold-shaped
data (D1 make\_moons, D5 Fashion-MNIST), where one global bandwidth fits poorly. Parameters are $k$
and the affinity (with $\gamma$ or the neighbour count). A dense RBF affinity costs $O(n^2d)$ time and
$O(n^2)$ memory, and a full eigendecomposition $O(n^3)$; the sparse $k$-NN graph with an iterative
eigensolver is much cheaper but still limits $n$ (Section~\ref{sec:runtime}). Spectral clustering
needs $k$, has no noise label, and is sensitive to the affinity scale.
\subsection{Agglomerative clustering}
\habegin
% Source: hoanganh/src/sc4020/clustering.py (wine_models, breast_cancer_models);
% ARI from hoanganh/results/{wine,breast_cancer}_results.csv.
\hatag{}Agglomerative clustering starts with every point as its own cluster and repeatedly merges the two
closest clusters until $k$ remain. The linkage defines which two are closest. We use Ward
linkage~\cite{ward1963}, which merges the pair whose union raises the total within-cluster variance
least. This is the K-means objective, searched greedily from the bottom up. The method therefore
needs no starting centres and returns the same partition on every run, which makes it a direct check
on K-means: same objective, different search. The only parameter is $k$, set to the number of classes
(3 on D6, 2 on D7); distances are Euclidean on standardised features. It needs $O(n^2)$ time and
memory for the pairwise distances, so we ran it on the two small tables only ($n=178$ and $n=569$).
Like K-means, it assumes compact clusters, assigns every point and has no noise label. Its own
weakness is that a merge is never undone, so an early mistake stays in the final partition. On both
datasets it matched the labels less well than K-means (ARI 0.790 against 0.897 on D6, 0.575 against
0.654 on D7).
\haend

\begin{table}[h]\centering\small
\caption{Structural assumptions of the methods compared.}
\label{tab:assump}
\resizebox{\textwidth}{!}{\begin{tabular}{lccccccc}
\toprule
& \textbf{K-Means} & \textbf{K-Means++} & \textbf{DBSCAN} & \textbf{HDBSCAN}
& \textbf{GMM} & \textbf{Spectral} & \hatag\textbf{Agglom.} \\
\midrule
Cluster shape assumed      & isotropic & isotropic & arbitrary & arbitrary & ellipsoidal & connected (graph) & compact (Ward) \\
Requires $k$               & yes & yes & no & no & yes & yes & yes \\
Explicit noise label       & no & no & yes & yes & no & no & no \\
Handles varying density    & no & no & no & yes & partly\textsuperscript{b} & partly\textsuperscript{c} & no \\
Parameter is interpretable & no & no & yes (km) & partly & no & no & no \\
Time complexity            & $O(nkdi)$ & $O(nkdi)$ & $O(n\log n)$\textsuperscript{a} & $O(n\log n)$\textsuperscript{a} & $O(nkd^2 i)$ & $O(n^2d+n^3)$\textsuperscript{d} & $O(n^2)$ \\
\bottomrule
\end{tabular}}
\par\raggedright\footnotesize\textsuperscript{a}\,With an effective spatial index; $O(n^2)$ worst case.
\textsuperscript{b}\,Per-component covariance and weight.
\textsuperscript{c}\,With a $k$-NN affinity, which adapts to local scale.
\textsuperscript{d}\,Dense RBF affinity and full eigendecomposition; lower with a sparse $k$-NN graph.
\end{table}
